% =========================================================
% Hasse Diagrams and the Duality Principle
% =========================================================

\section{Hasse Diagrams and the Duality Principle}\label{sec:hasse-duality}

% ---------------------------------------------------------
% TOOLKIT
% ---------------------------------------------------------
\begin{tcolorbox}[colback=gray!6, colframe=gray!40, arc=2pt,
  left=6pt, right=6pt, top=4pt, bottom=4pt,
  title={\small\textbf{Hasse Diagrams and Duality — Quick Reference}},
  fonttitle=\small\bfseries]
\small
\begin{tabular}{l l l}
\toprule
\textbf{Concept} & \textbf{Meaning} & \textbf{Detail} \\
\midrule
Cover relation    & $a \lessdot b$: $a < b$ with nothing strictly between       & \hyperref[def:cover-relation]{↓ Def} \\
Hasse diagram     & Graph encoding a poset; edges only for cover relations      & \hyperref[def:hasse]{↓ Def} \\
Dual poset        & $(A, \geq)$ obtained by reversing the order on $(A, \leq)$  & \hyperref[def:dual-poset]{↓ Def} \\
Duality principle & Every theorem about posets yields a dual theorem for free   & \hyperref[prop:order-duality]{↓ Prop} \\
\bottomrule
\end{tabular}
\end{tcolorbox}

\vspace{1em}

A partial order can be an enormous set of pairs, impossible to visualise
directly.  The \emph{Hasse diagram} is a compression: it retains only the
pairs that are ``immediate'' in the order, discarding those implied by
transitivity and reflexivity.  From the compressed diagram the full order
can be recovered by taking transitive closure.  Hasse diagrams are the
standard visual language for finite posets; knowing how to read and draw them
is essential for understanding the order examples in topology and lattice theory.

% ---------------------------------------------------------
% Cover Relation
% ---------------------------------------------------------

\begin{tcolorbox}[colback=propbox, colframe=propborder, arc=2pt,
  left=6pt, right=6pt, top=4pt, bottom=4pt,
  title={\small\textbf{Definition (Cover Relation)}},
  fonttitle=\small\bfseries]
\label{def:cover-relation}
Let $(A, \leq)$ be a poset and $a, b \in A$.  We say $b$ \emph{covers} $a$,
written $a \lessdot b$, if $a < b$ and there is no $c \in A$ with
$a < c < b$.
\end{tcolorbox}

\begin{remark*}[Fully quantified form]
$a \lessdot b \;\leftrightarrow\; (a < b) \;\wedge\;
\neg(\exists c \in A,\; a < c \;\wedge\; c < b)$.
\end{remark*}

\begin{remark}[Intuition]
\label{rem:cover-intuition}
The cover relation strips the partial order down to its ``essential edges.''
If $a \lessdot b$, then $b$ is the \emph{immediate successor} of $a$: you
cannot insert anything between them.  In a finite poset, every strict
relation $a < b$ factors as a finite chain of cover steps:
$a = a_0 \lessdot a_1 \lessdot \cdots \lessdot a_k = b$.  In an infinite
poset, this factorisation may fail (e.g., in $(\mathbb{R}, \leq)$, no two
distinct real numbers are in a cover relation because there is always a
rational strictly between them).
\end{remark}

% ---------------------------------------------------------
% Hasse Diagram
% ---------------------------------------------------------

\begin{tcolorbox}[colback=propbox, colframe=propborder, arc=2pt,
  left=6pt, right=6pt, top=4pt, bottom=4pt,
  title={\small\textbf{Definition (Hasse Diagram)}},
  fonttitle=\small\bfseries]
\label{def:hasse}
The \emph{Hasse diagram} of a finite poset $(A, \leq)$ is the directed graph
whose vertices are the elements of $A$, with an upward edge from $a$ to $b$
whenever $a \lessdot b$.  Three graphical conventions apply:
\begin{enumerate}[label=(\roman*)]
  \item Reflexive loops ($a \leq a$) are suppressed.
  \item Edges implied by transitivity are suppressed: if $a \lessdot b$ and
        $b \lessdot c$, no direct edge from $a$ to $c$ is drawn.
  \item Direction is encoded by height: $a \lessdot b$ is drawn with $b$
        strictly above $a$, so arrowheads are unnecessary.
\end{enumerate}
\end{tcolorbox}

\begin{remark}[How to recover the full order]
\label{rem:hasse-recovery}
The full order can be read back from the Hasse diagram: $a \leq b$ if and
only if either $a = b$ or there is an upward path in the diagram from $a$ to
$b$ (a sequence of cover edges going strictly upward).  The diagram is therefore
a lossless encoding of the partial order, compressed by removing the redundant
reflexive and transitive edges.
\end{remark}

\begin{example}[Divisibility on $\{1,2,3,4,6,12\}$]
\label{ex:hasse-divisibility}
Order the set $D = \{1, 2, 3, 4, 6, 12\}$ by divisibility ($a \leq b$ iff
$a \mid b$).  The covers are:
\[
  1 \lessdot 2,\quad
  1 \lessdot 3,\quad
  2 \lessdot 4,\quad
  2 \lessdot 6,\quad
  3 \lessdot 6,\quad
  4 \lessdot 12,\quad
  6 \lessdot 12.
\]
The Hasse diagram places $1$ at the bottom and $12$ at the top.  The relation
$1 \leq 4$ holds in the poset but the edge $1 \to 4$ is \emph{not} drawn
because $1 \lessdot 2 \lessdot 4$ provides an upward path.  Drawing the direct
edge would be redundant.

Notice also: $4$ and $6$ are incomparable, so there is no edge between them —
they sit at the same level with no path connecting them.  The join
$4 \vee 6 = 12 = \mathrm{lcm}(4,6)$ and the meet $4 \wedge 6 = 2 = \gcd(4,6)$
are readable from the diagram as the lowest common ancestor and the highest
common descendant, respectively.
\end{example}

\begin{example}[Power set $\mathcal{P}(\{a,b,c\})$ under inclusion]
\label{ex:hasse-power-set}
The Hasse diagram of $(\mathcal{P}(\{a,b,c\}), \subseteq)$ has four levels:
$\varnothing$ at the bottom; the three singletons $\{a\},\{b\},\{c\}$ one
level up; the three two-element sets $\{a,b\},\{a,c\},\{b,c\}$ above those;
and $\{a,b,c\}$ at the top.  Each singleton is covered by the two two-element
sets containing it.  This diagram is the Boolean lattice $B_3$: it has a
bottom ($\varnothing$), a top ($\{a,b,c\}$), and every element has a complement
($A' = \{a,b,c\} \setminus A$).
\end{example}

\begin{remark}[Limitation to finite posets]
\label{rem:hasse-infinite}
For infinite posets such as $(\mathbb{N}, \leq)$, the cover relation is
well-defined ($n \lessdot n+1$) but the full diagram cannot be drawn.  In
$(\mathbb{R}, \leq)$ or $(\mathbb{Q}, \leq)$, the cover relation is empty:
no two elements are in a cover relation.  Hasse diagrams are a practical
visualisation tool for finite or schematic examples and should not be used
as a proof device in infinite settings.
\end{remark}

% ---------------------------------------------------------
% Dual Poset
% ---------------------------------------------------------

\begin{tcolorbox}[colback=propbox, colframe=propborder, arc=2pt,
  left=6pt, right=6pt, top=4pt, bottom=4pt,
  title={\small\textbf{Definition (Dual Poset)}},
  fonttitle=\small\bfseries]
\label{def:dual-poset}
Let $(A, \leq)$ be a poset.  The \emph{dual poset} is $(A, \geq)$, where
$a \geq b$ is defined to mean $b \leq a$.  The dual is also written
$(A, \leq^{\mathrm{op}})$.
\end{tcolorbox}

\begin{remark}[Intuition]
\label{rem:dual-intuition}
The dual poset is obtained by flipping the Hasse diagram upside down.  Every
structural feature is preserved but reflected: elements at the top are now at
the bottom, minimal elements become maximal, upper bounds become lower bounds,
and suprema become infima and vice versa (as established in
Proposition~\ref{prop:sup-inf-duality}).
\end{remark}

% ---------------------------------------------------------
% Duality Principle
% ---------------------------------------------------------

\begin{proposition}[\texorpdfstring{\hyperref[prf:order-duality]{Duality Principle for Posets}}{Duality Principle for Posets}]
\label{prop:order-duality}
Let $\Phi$ be any first-order statement about a poset $(A, \leq)$ expressed
using only the relation $\leq$.  Let $\Phi^*$ be the statement obtained from
$\Phi$ by replacing every occurrence of $\leq$ with $\geq$ (equivalently,
working in the dual poset).  Then:
\[
  (A, \leq) \models \Phi
  \;\;\Longleftrightarrow\;\;
  (A, \geq) \models \Phi.
\]
In particular, if $\Phi$ is a theorem about all posets, then so is $\Phi^*$.
\end{proposition}

\begin{remark*}[Fully quantified form]
$\forall$ posets $(A,\leq)$, $\forall$ first-order sentences $\Phi$ in the
language of $\leq$: $(A,\leq)\models\Phi \;\leftrightarrow\; (A,\geq)\models\Phi^*$,
where $\Phi^*$ substitutes $\geq$ for every $\leq$.
\end{remark*}

\begin{remark}[Theorems come in pairs at no extra cost]
\label{rem:duality-pairs}
The duality principle means every proved theorem about posets automatically
yields a second theorem — its dual — without any additional work.  You will
encounter this repeatedly in analysis.  Here are the most important instances:
\begin{itemize}
  \item \emph{Supremum $\leftrightarrow$ Infimum.}  Every theorem about
        $\sup$ yields a theorem about $\inf$ by duality.  In particular,
        the completeness axiom for $\mathbb{R}$ (every nonempty bounded-above
        subset has a supremum) immediately implies that every nonempty
        bounded-below subset has an infimum — not a separate axiom but a
        consequence by duality.
  \item \emph{Minimal $\leftrightarrow$ Maximal.}  A statement about minimal
        elements dualises to a statement about maximal elements.  Zorn's Lemma
        (every poset in which every chain is bounded above has a maximal
        element) dualises to: every poset in which every chain is bounded below
        has a minimal element.
  \item \emph{Well-ordering is not self-dual.}
        $(\mathbb{N}, \leq)$ is well-ordered (every nonempty subset has a
        least element); its dual $(\mathbb{N}, \geq)$ is not well-ordered
        because the whole set has no greatest element.  Duality preserves
        the form of a definition but not whether a particular structure
        satisfies it.
  \item \emph{Join $\leftrightarrow$ Meet in lattices.}  Every lattice identity
        involving $\vee$ and $\wedge$ dualises by swapping the two operations.
        The distributive law $a \wedge (b \vee c) = (a \wedge b) \vee
        (a \wedge c)$ dualises to $a \vee (b \wedge c) =
        (a \vee b) \wedge (a \vee c)$: proving one gives the other for free.
\end{itemize}
\end{remark}

\begin{remark}[Transition]
\label{rem:hasse-transition}
The order vocabulary developed across the order sections — partial orders,
total orders, well-orders, bounds, suprema, infima, lattices, directed sets,
and duality — is the complete toolkit for order-theoretic reasoning in analysis.
In Volume~II, the construction of $\mathbb{R}$ requires only the completeness
axiom (existence of suprema) and the Dedekind cut or Cauchy sequence
constructions, both of which use only the vocabulary established here.  In
Volume~III, metric space arguments, the definition of compactness, and the
construction of Lebesgue measure all draw on this vocabulary repeatedly.
\end{remark}
